% !TeX root = ../main.tex

\section{Introduction}\label{sec:introduction}

Over recent years, social media influencers have emerged as powerful actors shaping consumer preferences and financial decision-making. Global spending on influencer marketing reached \$24~billion in 2024 \cite{pan2025}, reflecting their growing economic significance. Audiences often emulate their online opinion leaders, and influencer-driven content has been shown to significantly affect consumer behaviour~\cite{libai2025influencer}. This influence increasingly extends into financial markets, where digital communities can shape investment decisions. The case of Reddit's \texttt{WallStreetBets} illustrates this phenomenon, as collective attention drove waves of high-risk retail trades \cite{warkulat2024}.

In the blockchain and cryptocurrency sector, a distinct class of influencers, often termed \emph{Key Opinion Leaders (KOLs)} or ``crypto-influencers,'' has emerged. These individuals command large followings and play a central role in shaping sentiment and investment behaviour. Empirical research shows that crypto influencers' communications can drive short-term cryptocurrency price fluctuations and trading volumes, underscoring their market-moving capacity \cite{Merkley2024,Moser2023}. Unlike mainstream lifestyle influencers, crypto KOLs operate in highly volatile, loosely regulated markets characterised by pseudonymous participation, complex tokenomics, and irreversible transactions. They often hold positions in the assets they discuss, promote or build products around, and communicate with communities that are both financially and ideologically invested. In this setting, a single post about a low-cap token may move prices, expose followers to pump-and-dump schemes or rug pulls, and generate substantial losses for retail investors with limited technical or financial expertise. Credibility is therefore not merely an abstract reputational asset but a determinant of real financial risk.

Crypto KOLs differ from generic financial influencers (finfluencers) and traditional financial advisors. Whereas finfluencers frequently comment on diversified portfolios or macro trends, crypto counterparts position themselves as domain specialists in decentralised finance (DeFi) and Web3 protocols, translating highly technical developments into actionable narratives. They may communicate under pseudonyms while simultaneously leaving public traces on-chain (e.\,g., via visible wallet holdings or governance activity), creating new forms of transparency and conflict of interest. Regulatory guidance for such actors is still emerging and fragmented across jurisdictions, leaving many KOLs to self-interpret obligations around disclosure, suitability, and market integrity. These socio-technical conditions make questions of credibility, ethics, and responsibility especially acute.

Academic attention to finfluencers remains still limited. Prior work has examined how influencers build personal brands and motivations in mainstream contexts~\cite{torhonen2019}, how sponsorship disclosure and perceived authenticity affect persuasion~\cite{boerman2017}, and how finfluencers and celebrity figures can move crypto markets in the short run~\cite{Merkley2024,AnteMuskCrypto2023}. Parallel literatures in human-computer interaction (HCI) and usable security document how cryptocurrency users struggle with risks, scams, and security practices~\cite{Krombholz2016,Mai2020,SasKhairuddin2017,Abramova2021}. However, there is still a limited understanding about how crypto KOLs become and remain opinion leaders in such high-stakes environments, how they conceptualise and cultivate credibility, and how they navigate ethical tensions between education, promotion, and self-interest.

Guided by self-determination theory (SDT) \citep{deci2000} and using qualitative research methods, this study examines the motivations, beliefs, and communicative strategies of crypto KOLs. Specifically, we investigate:
\begin{itemize}
  \item \textbf{RQ1:} What extrinsic and intrinsic factors motivate individuals to become and remain key opinion leaders in the cryptocurrency ecosystem?
  \item \textbf{RQ2:} What defining characteristics and everyday practices set crypto KOLs apart from other influencers and members of the crypto community?
  \item \textbf{RQ3:} How do crypto KOLs conceptualise and enact credibility, ethics, and responsibility toward their communities in high-risk, speculative markets?
\end{itemize}

\noindent \textbf{Contributions.} This study advances understanding of crypto-influencer dynamics through four interrelated contributions which are grounded in theory, empirical evidence, methodological reflection, and practical impact.

\emph{Theoretical contribution.} We apply self-determination theory to high-risk financial contexts by showing how the needs for autonomy, competence, and relatedness are negotiated under conditions of volatility, sponsorship pressure, and shared financial exposure with followers. Our analysis surfaces tensions between autonomy and monetisation, between technical competence and communicative accessibility, and between relatedness and the risk of encouraging speculation. We further contribute to influencer marketing and creator-labour research by theorising crypto KOLs as hybrid actors who simultaneously serve as educators, informal financial advisors, entrepreneurs, and community stewards. Whereas prior work conceptualises credibility in terms of static credentials such as formal education, professional background, and financial certifications~\cite{NurhandayaniWhoDeserves2025}, our findings reveal credibility as a \emph{self-determined, ethically enacted practice} rooted in SDT needs and shaped by KOLs' own norms, constraints, and responsibilities. We identify four emergent dimensions that function as community-recognised markers of trustworthiness: (i)~\emph{self-regulation and voluntary constraint}, whereby KOLs decline misaligned sponsorships and impose personal rules on promotion; (ii)~\emph{bounded epistemic competence}, acknowledging the limits of one's expertise and avoiding prognostication; (iii)~\emph{accountability}, cultivating long-term trust through transparent disclosure and community stewardship; and (iv)~\emph{reflexive self-correction}, learning from past failures and continuously reassessing own practices. This reframes credibility as a socio-technical, self-regulated performance rather than a static credential.

\emph{Empirical contribution.} Drawing on thirteen interviews with crypto KOLs from Europe, the United States, and Asia (traders, educators, founders, and analysts), we present a thematic analysis that details (i) extrinsic drivers (e.\,g., sponsorship revenue, monetisation of analytics, social capital), (ii) intrinsic motives (e.\,g., enjoyment, educating newcomers, ideological commitment, desire for mastery and community), (iii) characteristics and practices that distinguish crypto KOLs from mainstream influencers (e.\,g., technical expertise, regulatory literacy, use of on-chain signals), and (iv) ethical and community norms that guide their communication strategies (e.\,g., disclosure practices, self-imposed restrictions on promotions). These qualitative insights fill a notable gap in the literature on digital opinion leadership by providing rich, contextualised accounts of a previously understudied influencer category.

\emph{Methodological contribution.} To systematically analyse a complex and rapidly evolving domain, we describe and critically reflect on a hybrid workflow that couples conventional qualitative interviewing with large language model--assisted thematic analysis. In our approach, a large language model (LLM) proposes candidate codes and SDT mappings on anonymised transcript segments, while human researchers retain full control over the codebook development, theme construction, and interpretation. We show how this human-in-the-loop workflow can broaden candidate theme coverage and support transparency without treating the LLM as an autonomous analyst or source of ground truth.

\emph{Practical contribution.}
Our findings yield actionable recommendations for industry and policy. For marketers, campaigns that resonate with KOLs' community-oriented motives such as education and protection are more likely to produce credible and enduring partnerships than those based solely on one-off paid promotions. For platform designers, we suggest features such as transparent disclosure badges, longitudinal accuracy or ``track record'' indicators, and cross-platform reputation profiles that help users identify credible actors and incentivise responsible behaviour. For regulators, the study highlights the need for proportionate disclosure rules, registration thresholds for high-reach influencers, and cross-jurisdictional coordination to address scams, rug pulls, and other forms of market manipulation while preserving the participatory ethos of Web3. Together, these measures offer a roadmap for strengthening trustworthy influence in cryptocurrency and for designing socio-technical infrastructures that support ethical, evidence-based content.

\noindent \textbf{Paper structure.} The remainder of this paper is organised as follows. Section~\ref{sec:background} reviews relevant literature on cryptocurrency user behaviour, social media influencers, finfluencers, creator labour, and the SDT framework. Section~\ref{sec:methodology} describes our qualitative methodology, including interview procedures, participant demographics, and coding strategy. Section~\ref{sec:results} presents the main results structured along the specified RQs. In Section~\ref{sec:discussion}, we discuss the implications of these findings, limitations and future research directions, and Section~\ref{sec:conclusion} concludes with a summary of contributions.
